\documentclass[10pt,a4paper]{amsart}
\usepackage[margin=19mm]{geometry}
\usepackage{amsmath,amssymb,mathtools}
\usepackage[scr=rsfs,cal=euler]{mathalfa}
\usepackage{xfrac}
\usepackage[unicode,hidelinks,bookmarks=false]{hyperref}
\newcommand{\ie}{i.e.\ }
\newcommand{\cf}{cf.\ }
\newcommand{\resp}{respectively\ }
\newcommand{\Subsection}{\subsection}
\providecommand{\nobreakdash}{\nobreak\hbox{-}}
\setlength{\emergencystretch}{2em}
\multlinegap0pt
\usepackage{bm}
\usepackage{bbm}
\usepackage{dsfont}
\usepackage{xspace}
\usepackage{cancel}
\usepackage{mathtools}
\usepackage{cleveref}
\usepackage{amsthm}
\makeatletter
\@ifundefined{theorem}{\newtheorem{theorem}{Theorem}}{}
\@ifundefined{claim}{\newtheorem{claim}[theorem]{Claim}}{}
\@ifundefined{fact}{\newtheorem{fact}[theorem]{Fact}}{}
\@ifundefined{lemma}{\newtheorem{lemma}[theorem]{Lemma}}{}
\makeatother
\makeatletter
\newcommand{\Boo}{\{0,1 \}}
\newcommand{\F}{\mathbb{F}}
\newcommand{\bigO}{\mathcal{O}}
\newcommand{\brac}[1]{\left[ #1 \right]}
\newcommand{\ml}{\mathsf{ml}}
\expandafter\ifx\csname myq\endcsname\relax
\newenvironment{myq}{\newline\noindent\begin{minipage}{\textwidth}}{\end{minipage}}
\fi
\newcommand{\paren}[1]{\left( #1 \right)}
\newcommand{\set}[1]{\left\{ #1 \right\}}
\makeatother
\title[New Bounds for the Ideal Proof System in Positive Characteri]{New Bounds for the Ideal Proof System in Positive Characteristic}
\author{Amik Raj Behera, Nutan Limaye, Varun Ramanathan, Srikanth Srinivasan}
\date{Source: arXiv:2506.16397v1 (CC BY 4.0)}
\begin{document}
\maketitle

\noindent\emph{Source and licence.} New Bounds for the Ideal Proof System in Positive Characteristic. Version arXiv:2506.16397v1 (CC BY 4.0), distributed under the Creative Commons Attribution 4.0 International licence, as recorded in the arXiv licence metadata for this exact version and shown on the abstract page https://arxiv.org/abs/2506.16397v1. Full rights notice: licenses/arxiv-2506.16397-CC-BY-4.0.txt.\par\smallskip

\noindent\textit{Editorial scope.} Counted unit: the Claim whose source label is \texttt{claim:multilinearize-exponent-vector} in the article (source0.tex lines 1881--1894, source label \texttt{claim:multilinearize-exponent-vector}), delivered together with the article's own complete original proof of it (source0.tex lines 1895--1948, one \texttt{proof} environment of 54 lines). No mathematics is written, repaired or re-derived: statement and proof are the article's own text line for line. Required context retained uncounted: fact:multilinear (lines 756--763); lemma:ml-prod-two-elem (lines 1820--1833). Every definition, display and label of those lines is retained unchanged. Omitted from this package and not redistributed: 1--755, 764--1819, 1834--1880, 1949--2229. Nothing inside a retained excerpt is omitted or altered: every retained body is delivered exactly as the article's lines are written, so no omission receipt of any kind accompanies this package. Citations: the retained excerpts carry no citation command at all, so no bibliography entry of the article is transcribed and no reference is redistributed. Reference adaptation: none was needed and none was made; every reference inside the delivered unit resolves inside it, so no unresolved reference or citation is produced. Printed slips kept literal: no printed word, symbol, hypothesis or number of the counted statement or of its proof was altered. Layout and class: the article's own class is \texttt{article} with packages including csquotes, geometry, fontenc, inputenc, authblk, amsmath, amssymb, amsthm, thmtools, thm-restate, mathabx, fullpage, mdwlist, graphicx; the wrapper is the lane's pinned \texttt{amsart} editorial wrapper at 10pt on A4, which restates the theorem-like environments the retained text needs and adds the article's own macro definitions that the retained text needs (\texttt{\textbackslash Boo}, \texttt{\textbackslash F}, \texttt{\textbackslash bigO}, \texttt{\textbackslash brac}, \texttt{\textbackslash ml}, \texttt{\textbackslash paren}, \texttt{\textbackslash set}), with extra wrapper packages (bm, bbm, dsfont, xspace, cancel, mathtools, cleveref). The delivered numbering is the unit's own. Assets: no retained span contains a figure environment, a graphics-inclusion command, a PDF-page inclusion, a table or a TikZ picture; no image file is copied into this package, the package declares no asset dependency and no image file is redistributed with it. Licence: version arXiv:2506.16397v1 of this article is distributed under the Creative Commons Attribution 4.0 International licence, as recorded in the arXiv licence metadata for this exact version and shown on the abstract page https://arxiv.org/abs/2506.16397v1. A full rights notice is delivered at licenses/arxiv-2506.16397-CC-BY-4.0.txt. Characters: every retained excerpt is pure ASCII, so the delivered engine, XeLaTeX with its default Latin Modern text font, reports no missing character.
\begin{fact}[Standard facts on multilinear polynomials]\label{fact:multilinear}
Let $f(\mathbf{x}), g(\mathbf{x}) \in \F[\mathbf{x}]$.
\begin{itemize}
    \item $f(\mathbf{x})$ and $\ml[f(\mathbf{x})]$ agree on the Boolean cube $\Boo^{n}$.
    \item $f(\mathbf{x})$ and $g(\mathbf{x})$ agree on the Boolean cube $\Boo^{n}$ if and only if $\ml[f(\mathbf{x})]$ is equal to the $\ml[g(\mathbf{x})]$.
    \item $\ml[f(\mathbf{x}) g(\mathbf{x})] = \ml[\ml[f(\mathbf{x})] \ml[g(\mathbf{x})]]$.
\end{itemize}
\end{fact}

\begin{lemma}[Multilinearization of product of two elementary symmetric polynomials]\label{lemma:ml-prod-two-elem}
Fix any two natural numbers $\alpha$ and $\beta$. Then
\begin{itemize}
    \item There exists polynomials $R_{\alpha, \beta,j}(\mathbf{x})$'s such that
    \begin{align*}
        \ml[e_{\alpha}(\mathbf{x}) \cdot e_{\beta}(\mathbf{x})] \; = \; e_{\alpha}(\mathbf{x}) \cdot e_{\beta}(\mathbf{x}) - \sum_{j=1}^{n} R_{\alpha, \beta, j}(\mathbf{x}) \cdot (x_{j}^{2} - x_{j}),
    \end{align*}
    where each polynomial $R_{\alpha, \beta, j}(\mathbf{x})$ has a circuit of size $\bigO(n^{3})$ and depth $5$ (a $\Sigma \Pi \Sigma \Pi \Sigma$ circuit).
    \item There exists coefficients $c_{\alpha, \beta}^{(i)}$'s such that
    \begin{align*}
        \ml[e_{\alpha}(\mathbf{x}) \cdot e_{\beta}(\mathbf{x})] \; = \; \sum_{i = 1}^{n} c_{\alpha,\beta}^{(i)} \; e_{i}(\mathbf{x})
    \end{align*}
\end{itemize}
\end{lemma}

\begin{claim}\label{claim:multilinearize-exponent-vector}
Let $\bm{\mu} \in \set{0,1,\ldots,p-1}^{r}$ denote an exponent vector as described above. Then for any $k \in [|\bm{\mu}|]$, the following holds:
\begin{itemize}
    \item (Constant-depth circuit). There exists polynomials $R_{\leq \alpha_{k}, j}(\mathbf{x})$'s such that
    \begin{align*}
        \ml\brac{\prod_{\ell = 1}^{k} e_{\alpha_{\ell}}(\mathbf{x})} \; = \; \prod_{\ell=1}^{k} e_{\alpha_{\ell}}(\mathbf{x}) - \sum_{j=1}^{n} R_{\leq \alpha_{k}, j}(\mathbf{x}) \cdot (x_{j}^{2} - x_{j}),
    \end{align*}
    and the polynomials $R_{\leq \alpha_{k}, j}(\mathbf{x})$'s have circuits of size $\bigO(n^{3} \cdot k)$ and depth $5$.
    \item (Structure). There exists coefficients $c_{\leq \alpha_{k}}^{(1)}, \ldots, c_{\leq \alpha_{k}}^{(n)}$ such that
    \begin{align*}
        \ml\brac{\prod_{\ell = 1}^{k} e_{\alpha_{\ell}}(\mathbf{x})} \; = \; \sum_{i=1}^{n} c_{\leq \alpha_{k}}^{(i)} \; e_{i}(\mathbf{x})
    \end{align*}
\end{itemize}
\end{claim}

\begin{proof}[Proof of \Cref{claim:multilinearize-exponent-vector}]
We will prove this using induction on $k$.
\paragraph{Base case:}For $k = 1$, we have $\ml[e_{\alpha_{1}}] = e_{\alpha_{1}}(\mathbf{x})$ and $R_{\leq \alpha_{1}, j}(\mathbf{x}) = 0$. The claim holds for the base case.

\paragraph{Induction step:}Now we assume the induction is true for $k$ and prove it for $(k+1)$. We will first prove the \emph{constant-depth circuit} item for $(k+1)$ and then prove the \emph{structure} item for $(k+1)$.\\

\noindent
Using the \textbf{structure} item of the induction hypothesis, we have,
\begin{equation}\label{eqn:ml-char-p-1}
    \ml\brac{\prod_{\ell = 1}^{k} e_{\alpha_{\ell}}(\mathbf{x})} \; = \; \sum_{i=1}^{n} c_{\leq \alpha_{k}}^{(i)} \; e_{i}(\mathbf{x})
\end{equation}
Using the third item of \Cref{fact:multilinear},
\begin{align*}
    \ml\brac{\prod_{\ell = 1}^{k+1} e_{\alpha_{\ell}}(\mathbf{x})} \; = \; \ml\brac{\ml\brac{\prod_{\ell = 1}^{k} e_{\alpha_{\ell}}(\mathbf{x})} \cdot e_{\alpha_{k+1}}(\mathbf{x})} \\ \\
    = \ml\brac{\sum_{i = 1}^{n} c_{\leq \alpha_{k}}^{(i)} \; e_{i}(\mathbf{x}) \cdot e_{\alpha_{k+1}}(\mathbf{x}) } && (\text{Using \Cref{eqn:ml-char-p-1}}) 
\end{align*}
\begin{equation}\label{eqn:partial-ml-8}
    = \sum_{i = 1}^{n} c_{\leq \alpha_{k}}^{(i)} \; \ml[e_{i}(\mathbf{x}) \cdot e_{\alpha_{k+1}}(\mathbf{x}) ]
\end{equation}
For each $i \in [n]$, we apply the \textbf{constant-depth circuit} item from \Cref{lemma:ml-prod-two-elem} on $e_{i}(\mathbf{x}) \cdot e_{\alpha_{k+1}}(\mathbf{x})$ to get:
\begin{align*}
    \ml[e_{i}(\mathbf{x}) \cdot e_{\alpha_{k+1}}(\mathbf{x})] \; = \; e_{i}(\mathbf{x}) \cdot e_{\alpha_{k+1}}(\mathbf{x}) - \sum_{j=1}^{n} D_{i,\alpha_{k+1}}(\mathbf{x}) \cdot (x_{j}^{2} - x_{j}),
\end{align*}
where the polynomials $D_{i,\alpha_{k+1}}(\mathbf{x})$ have a circuit of size $\bigO(n^{3})$ and depth $5$. Substituting it in \Cref{eqn:partial-ml-8},
\begin{align*}
    \ml\brac{\prod_{\ell = 1}^{k+1} e_{\alpha_{\ell}}(\mathbf{x})} \; = \; \sum_{i=1}^{n} c_{\leq \alpha_{k}}^{(i)} \; \paren{e_{i}(\mathbf{x}) \cdot e_{\alpha_{k+1}}(\mathbf{x}) - \sum_{j=1}^{n} D_{i,\alpha_{k+1}}(\mathbf{x}) \cdot (x_{j}^{2} - x_{j})} \\ \\
    = \paren{\sum_{i = 1}^{n} c_{\leq \alpha_{k}}^{(i)} e_{i}(\mathbf{x}) } e_{\alpha_{k+1}}(\mathbf{x}) - \sum_{j=1}^{n} \paren{\sum_{i = 1}^{n} c_{\leq \alpha_{k}}^{(i)}  D_{i,\alpha_{k+1}}(\mathbf{x})} \; \cdot \; (x_{j}^{2} - x_{j})
\end{align*}
Using \Cref{eqn:ml-char-p-1},
\begin{align*}
    \sum_{i = 1}^{n} c_{\leq \alpha_{k}}^{(i)} e_{i}(\mathbf{x}) \; = \; \ml\brac{\prod_{\ell = 1}^{k} e_{\alpha_{\ell}}(\mathbf{x})} \; = \; \prod_{\ell = 1}^{k} e_{\alpha_{\ell}}(\mathbf{x}) - \sum_{j=1}^{n} R_{\leq \alpha_{k}, j}(\mathbf{x}) \cdot (x_{j}^{2} - x_{j}),
\end{align*}
where we use the \textbf{constant-depth circuit} item from the induction hypothesis for the last equality. The polynomials $R_{\leq \alpha_{k}, j}(\mathbf{x})$'s have circuits of size $\bigO(n^{3} \cdot k)$ and depth $5$ (a $\Sigma \Pi \Sigma \Pi \Sigma$ circuit). Using this in the previous expression, we have
\begin{gather*}
    \ml\brac{\prod_{\ell = 1}^{k+1} e_{\alpha_{\ell}}(\mathbf{x})} \\ 
    = \; \paren{ \prod_{\ell = 1}^{k} e_{\alpha_{\ell}}(\mathbf{x}) - \sum_{j=1}^{n} R_{\leq \alpha_{k}, j}(\mathbf{x}) \cdot (x_{j}^{2} - x_{j}) } e_{\alpha_{k+1}}(\mathbf{x}) - \sum_{j=1}^{n} \paren{\sum_{i = 1}^{n} c_{\leq \alpha_{k}}^{(i)}  D_{i,\alpha_{k+1}}(\mathbf{x})} \; \cdot \; (x_{j}^{2} - x_{j}) \\
    = \; \prod_{\ell=1}^{k+1} e_{\alpha_{\ell}}(\mathbf{x}) - \sum_{j=1}^{n} \; \underbrace{\paren{\sum_{i = 1}^{n} c_{\leq \alpha_{k}}^{(i)}  D_{i,\alpha_{k+1}}(\mathbf{x}) + R_{\leq \alpha_{k}, j}(\mathbf{x})}}_{:= R_{\leq \alpha_{k+1}, j}(\mathbf{x})} \; \cdot \; (x_{j}^{2} - x_{j})
\end{gather*}
The polynomial $R_{\leq \alpha_{k+1},j}(\mathbf{x})$ has a circuit of size $\bigO(n^{3} \cdot (k+1))$ and depth $5$ (a $\Sigma \Pi \Sigma \Pi \Sigma$ circuit). This shows the \emph{constant-depth circuit} item of the induction.\\


\noindent
By applying the \textbf{structure item} of \Cref{lemma:ml-prod-two-elem} on $e_{i}(\mathbf{x}) \cdot e_{\alpha_{k+1}}(\mathbf{x})$, we get,
\begin{align*}
     \ml\brac{\prod_{\ell = 1}^{k} e_{\alpha_{\ell}}(\mathbf{x})} \; = \; \sum_{i=1}^{n} c_{\leq \alpha_{k}}^{(i)} \; e_{i}(\mathbf{x})
\end{align*}
Substituting it in \Cref{eqn:partial-ml-8},

\begin{align*}
    \ml\brac{\prod_{\ell = 1}^{k+1} e_{\alpha_{\ell}}(\mathbf{x})} \; = \; \sum_{i = 1}^{n} c_{\leq \alpha_{k}}^{(i)} \; \sum_{i'=1}^{n} d_{i}^{(i')} e_{i'}(\mathbf{x}) \; = \; \sum_{i = 1}^{n} c_{\leq \alpha_{k+1}}^{(i)} \; e_{i}(\mathbf{x}),
\end{align*}
where $c_{\leq \alpha_{k+1}}^{(i)} = \sum_{j=1}^{n} c_{\leq \alpha_{k}}^{(j)} d_{j}^{(i)}$. This completes the \emph{structure} item of the induction.\newline
This finishes the induction and thus we have finished the proof of \Cref{claim:multilinearize-exponent-vector}.
\end{proof}

\end{document}
